\apartado{2.4}{Teoría combinatoria}

\noindent\textit{Contar sin enumerar. La única pregunta que hay que resolver
antes de elegir la fórmula es: ¿importa el orden?}
\recordar{%
  $n! = n\cdot(n-1)\cdots 2\cdot 1$, y $0!=1$.\quad
  \textbf{Permutación} (importa el orden):
  $P(n,k)=\dfrac{n!}{(n-k)!}$.\quad
  \textbf{Combinación} (no importa):
  $C(n,k)=\dbinom{n}{k}=\dfrac{n!}{k!\,(n-k)!}$.\\[0.3em]
  Siempre $C(n,k) \le P(n,k)$: agrupar los órdenes reduce la cuenta.}

\ejercicio{0}{%
  Calculá $4!$ y $0!$. Escribí el desarrollo de $4!$, no sólo el resultado.}
\espacioresolver[1.6cm]
\respuesta{$4! = 24$; $0! = 1$.}
\solucion{%
  $4! = 4\times 3\times 2\times 1 = 24$.
  \\[0.3em]
  $0! = 1$ por definición. Suena arbitrario, pero es lo que hace que las
  fórmulas funcionen: hay exactamente \emph{una} manera de no elegir nada.}

\ejercicio{1}{%
  Tenés 3 fichas marcadas A, B y C. ¿De cuántas maneras distintas podés
  ordenarlas en fila? Enumeralas todas y comprobá que coincide con $3!$.}
\espacioresolver[2.4cm]
\respuesta{$3!=6$: ABC, ACB, BAC, BCA, CAB, CBA.}
\solucion{%
  $3! = 3\times 2\times 1 = 6$, y son:
  ABC, ACB, BAC, BCA, CAB, CBA.
  \\[0.3em]
  El razonamiento detrás del factorial: para el primer lugar hay 3 opciones,
  para el segundo quedan 2, para el tercero 1. $3\times 2\times 1$.}

\ejercicio{1}{%
  En un grupo de 5 estudiantes hay que formar parejas para un trabajo.
  ¿Cuántas parejas distintas se pueden formar?}
\espacioresolver[2cm]
\respuesta{$C(5,2)=10$.}
\solucion{%
  El orden no importa: la pareja «Ana y Beto» es la misma que «Beto y Ana».
  Entonces es combinación:
  \[ C(5,2) = \frac{5!}{2!\,3!} = \frac{120}{2\times 6} = 10 \]}

\ejercicio{2}{%
  De 5 estudiantes hay que elegir 3 para tres roles distintos: quien expone,
  quien toma notas y quien coordina. ¿Cuántas asignaciones distintas hay?}
\espacioresolver[2.2cm]
\respuesta{$P(5,3)=60$.}
\solucion{%
  Acá el orden \textbf{sí} importa: no es lo mismo que Ana exponga y Beto
  coordine que al revés. Es permutación:
  \[ P(5,3) = \frac{5!}{(5-3)!} = \frac{120}{2} = 60 \]
  O directamente: $5 \times 4 \times 3 = 60$.}

\ejercicio{2}{%
  Calculá $C(6,2)$ y $P(6,2)$. ¿Cuántas veces más grande es la permutación?
  Explicá por qué es exactamente ese factor.}
\espacioresolver[2.6cm]
\respuesta{$C(6,2)=15$, $P(6,2)=30$. La permutación es $2!=2$ veces más grande.}
\solucion{%
  \[ C(6,2)=\frac{6!}{2!\,4!}=15 \qquad P(6,2)=\frac{6!}{4!}=30 \]
  La permutación es exactamente $2$ veces más grande, y ese 2 es $2!$: cada
  grupo de 2 personas se puede ordenar de $2!=2$ maneras.
  \\[0.3em]
  En general $P(n,k) = C(n,k)\times k!$. Ésa es toda la diferencia entre las
  dos fórmulas, y saberlo evita memorizarlas por separado.}

\ejercicio{3}{%
  De los \textbf{43 estudiantes} que dieron positivo en el tamizaje, el
  servicio quiere elegir a 5 para una entrevista de profundización. El orden
  en que se los elige no importa. ¿De cuántas maneras distintas puede
  formarse ese grupo?}
\espacioresolver[3cm]
\respuesta{$C(43,5) = 962\,598$.}
\solucion{%
  El orden no importa: es un grupo, no una fila.
  \[ C(43,5) = \frac{43!}{5!\,38!}
     = \frac{43\times 42\times 41\times 40\times 39}{5\times 4\times 3\times 2\times 1}
     = \frac{115\,511\,760}{120} = 962\,598 \]
  \\[0.3em]
  \textbf{Truco de cálculo}: no calcules $43!$ — es un número astronómico.
  Alcanza con multiplicar los 5 factores de arriba y dividir por $5!$.}

\ejercicio{3}{%
  El mismo servicio ahora quiere \textbf{agendar} a 5 de esos 43 estudiantes
  en cinco horarios distintos del lunes. ¿De cuántas maneras puede hacerlo?
  Compará con el resultado anterior y explicá la diferencia.}
\espacioresolver[3cm]
\respuesta{$P(43,5) = 115\,511\,760$, que es $5! = 120$ veces mayor.}
\solucion{%
  Ahora el orden importa, porque cada persona queda en un horario distinto:
  \[ P(43,5) = 43\times 42\times 41\times 40\times 39 = 115\,511\,760 \]
  \\[0.3em]
  Es exactamente $120$ veces el resultado anterior, y $120 = 5!$: cada grupo
  de 5 personas se puede repartir en los cinco horarios de $5!$ maneras.
  \\[0.3em]
  \textbf{La misma gente, dos preguntas distintas.} Elegir a quiénes entrevistar
  es una combinación; decidir además cuándo, una permutación.}

\ejercicio{4}{%
  Para cada situación, decidí si hay que usar combinación o permutación, y
  calculá el resultado. Justificá la elección en una línea.
  \begin{enumerate}[label=(\alph*),topsep=2pt,itemsep=1pt]
    \item Elegir 3 de los 43 positivos para un grupo focal.
    \item Elegir 3 de los 43 y asignarles los roles de moderador, relator y observador.
    \item Elegir un comité de 2 personas entre 6 docentes.
    \item Elegir presidente y secretario entre 6 docentes.
  \end{enumerate}}
\espacioresolver[5cm]
\respuesta{(a) $C(43,3)=12\,341$; (b) $P(43,3)=74\,046$;
(c) $C(6,2)=15$; (d) $P(6,2)=30$.}
\solucion{%
  (a) \textbf{Combinación}: en un grupo focal todos cumplen el mismo papel.
  $C(43,3)=\dfrac{43\times42\times41}{6}=12\,341$.
  \\[0.3em]
  (b) \textbf{Permutación}: los roles son distintos.
  $P(43,3)=43\times42\times41=74\,046$.
  \\[0.3em]
  (c) \textbf{Combinación}: los dos miembros del comité son intercambiables.
  $C(6,2)=15$.
  \\[0.3em]
  (d) \textbf{Permutación}: presidente y secretario no son lo mismo.
  $P(6,2)=30$.
  \\[0.3em]
  \textbf{La pregunta que resuelve todo}: si intercambio dos de los elegidos,
  ¿cambia el resultado? Si cambia, importa el orden (permutación). Si no
  cambia, no importa (combinación). Nada más que eso.}
